\paragraph{What the EEG is for.} Ten second decisions from eight electrodes around the ears are right a little under two times in three, and one minute decisions nearly four times in five. That is too weak to steer audio by and strong enough to know, over a minute or two, whom a person keeps following. The architecture this implies puts head direction first, because it is instant, and uses the EEG as the slow estimate that keeps the chosen talker when the head turns away. The replay of real decisions under a wandering head says where that estimate earns its place, and where it does not: when a listener looks down at nothing, holding the last talker faced does as well without any EEG; when the listener looks at someone other than the person followed, or attention moves while the head stays, only the EEG keeps or finds the right voice. It also says the phone's fusion trusts a glance too much. Weighing the head half as much lets the belief hold through glances and catch most moves of attention the head does not make, while a real turn still comes through in about ten seconds; a head weight that grows with how long a talker is faced holds a little better and turns slower. The belief turns the weak decisions into that estimate, and the switch times measured on real EEG cut between listeners say how long it takes to let go of the old talker.

\paragraph{Calibration free, and held out.} Every number here is on people the decoder never saw, scaled only by their own first two minutes; the learned decoder is also held out by story content and by study, and an out of step control falls to chance throughout. The check that matters most for this data set is the one by content. Every listener in ds004015 heard the same two stories, so a network tested on new people could score well by recognising the stories rather than by reading attention. Trained without any block of the story content it was tested on, from anyone, the learned decoder lost nothing at one minute (\heldContentSixty{} against \heldStandardSixty{} with the content shared) and a few points at five minutes (\heldContentFiveMin{} against \heldStandardFiveMin{}); trained on one source study and tested on the other, it kept \heldStudySixty{} and \heldStudyFiveMin{}. These checks used one network per fold with the settings of the main search (Section~\ref{sec:decoders}), so they speak to what the decoder learns, not to its final accuracy. Set against the published numbers for electrodes around the ear (Table~\ref{tab:lit}), all of them from decoders trained on the person tested, a calibration free decoder over four electrodes per ear, each ear referenced to itself, is at least as accurate.


\paragraph{The contacts.} Calibration free, with each ear referenced to itself as two earbuds must be, dry contacts in the ear canal and concha carried no attention that a decoder trained on other people could read: \earInNn{} of one minute decisions with the learned decoder and \earInLin{} with the linear one, at a coin flip. Wet contacts around the same ears, at the same moments, gave \earAroundNn{} (\earAroundNnFive{} at five minutes), and putting the two sets together gave \earBothNn{}, less than around the ear alone: for people the decoder has never seen, the in-ear channels are noise to fit. The data set's own analysis, trained on each person, put dry in-ear decoding at 61.1\% and found that dropping the calibration cost it 11.7 points~\citep{geirnaert2025}, which is where this calibration free result sits. Two confounds remain, and the data cannot separate them: the in-ear contacts were dry and the around-ear ones wet, and the in-ear ones were fewer and closer together. Neither changes what an earbud that has to work out of the box needs, which is contacts that reach outside the canal. The specification that accompanies this work now requires at least two of the four sensing contacts per ear outside it.


\paragraph{The electronics.} The amplifier is not the limit. The noise, resolution and sampling rate an ADS1299 class part gives sit far inside what the decoder tolerates; so do the clock alignment the protocol achieves between two earbuds (about 10~ms, against \hwSkewHolds~ms that cost nothing) and the frame loss a Bluetooth link allows (the specification asks for at most 1\%; scattered loss of \hwLossHolds{} of frames cost under a point). What the decoder does not tolerate are speech envelopes misaligned with the EEG by a few tens of milliseconds and movement jolts left in the signal. Both are requirements on timing and software rather than on parts: the microphones' envelopes have to be stamped on the same clock as the electrodes, to about 10~ms, and the phone's chain gates movement out, as it now does.


\paragraph{The silent command.} The clench detector caught \clAround{} of cued clenches around the ear and \clHeadband{} on a consumer headband at its shipped setting, bimodally: those it missed clench too weakly at these electrodes for a threshold of four times the background, and lowering the threshold catches them at a price in false reports (\clAroundThree{} caught at \clAroundThreeFalse{} false an hour, against \clAroundFourFalse{} at the shipped setting). What sets it off is specific: swallowing (\clSwallow{} of swallows) and raised eyebrows (\clBrows{}), not speech (\clTongue{}). And a double clench almost never happens by itself: \clListenDouble{} an hour in \clListenHours{} hours of quiet listening, against \clListenOthers{} for singles and holds. A product should therefore put the command that must not fire by accident on the double clench, and reject swallows with a signal muscle activity does not give: the microphone that faces the ear canal in a noise cancelling earbud.
